\section{Related Works}
% Plain Language Summarization (PLS) is a popular area of study~\cite{Guo2020AutomatedLL,Goldsack2022MakingSS,Zaman2020HTSSAN,ji-etal-2024-rag,flores-etal-2023-medical}. The BioLaySum PLS shared task has been hosted by the BioNLP workshop multiple times, attracting more participants each year~\cite{chandrasekaran-etal-2020-overview-insights,biolaysumm-2023-overview,goldsack-etal-2024-biolaysumm}. 
PLS has typically treated readability as a binary feature - plain or technical~\cite{Guo2020AutomatedLL,Goldsack2022MakingSS,Zaman2020HTSSAN}. However, recent work has shown that preferred readability levels vary across different backgrounds~\cite{August2024KnowYA}. Current methods are limited, highlighting the need for additional research~\cite{luo-etal-2022-readability}. 
Past work in controllable summarization has focused primarily on attributes other than readability, such as topic and aspect~\cite{urlana-etal-2024-controllable, Zhang2022MACSumCS, he-etal-2022-ctrlsum} and on the news domain, rather than scientific~\cite{retkowski-waibel-2025-zero, Chan2021ControllableSW}. Finally, prior work in readability controllable summarization relies on compute and data intensive methods, such as reinforcement learning~\cite{ribeiro-etal-2023-generating}. 

To address this, we use control vectors, a representation-engineering method that manipulates hidden states during inference.  \textcite{zou2025representationengineeringtopdownapproach} introduced it for steering concepts like honesty. Several studies have identified limitations in control vectors \parencite{bartoszcze2025representationengineeringlargelanguagemodels, wehner2025taxonomyopportunitieschallengesrepresentation, vonrütte2024languagemodelsguidelatent, korznikov2025roguescalpelactivationsteering, stickland2024steeringeffectsimprovingpostdeployment}. Only a few works \parencite{10.5555/3737916.3742333, braun2025understandingunreliabilitysteeringvectors} have analyzed the underlying causes. None have proposed systematic methods for improving data quality, a gap our work addresses through our framework.


